\section{随机变量的多种收敛}
\label{sec:05-Probability/08-Limit-Theorems/01}

训练跑到第 40 轮，损失曲线看上去压平了。你在群里发一句：``收敛了。''

三个人追上来问，问的却是三件不同的事。做基础设施的想知道换一个随机种子重跑，这条曲线还会不会落在同一个地方；做评测的想知道误差平方的平均现在降到了多少；做统计的想知道把整套流程重复一百遍，那一百个终点会摆成什么形状。

三个问题的答案可以不一致。同一列随机变量，可能在一种意义下收敛，而在另一种意义下根本不收敛。概率论里的``收敛''不是一个词，是一组词，彼此强弱不同——这不是术语洁癖，因为下一节的弱大数定律与强大数定律，差的正是这组词里的一格。

\subsection{四种收敛，三条单向箭头}

先把关系摆出来，再解释每个名字。

\begin{center}
\begin{tikzpicture}[x=1cm,y=1cm,font=\small,>=stealth]
  \draw[black!45,fill=blue!10] (-1.9,0.8) rectangle (1.9,1.6);
  \node at (0,1.2) {几乎必然收敛};
  \draw[black!45,fill=red!8] (-1.9,-1.6) rectangle (1.9,-0.8);
  \node at (0,-1.2) {\(L^p\) 收敛};
  \draw[black!45,fill=green!12] (3.9,-0.4) rectangle (7.1,0.4);
  \node at (5.5,0) {依概率收敛};
  \draw[black!45,fill=orange!14] (8.9,-0.4) rectangle (12.1,0.4);
  \node at (10.5,0) {依分布收敛};
  \draw[->,very thick] (1.95,1.05) -- (3.85,0.3);
  \draw[->,very thick] (1.95,-1.05) -- (3.85,-0.3);
  \draw[->,very thick] (7.15,0) -- (8.85,0);
  \draw[dashed,black!55] (0,0.75) -- (0,0.28);
  \draw[dashed,black!55] (0,-0.28) -- (0,-0.75);
  \node[font=\footnotesize,black!70] at (0,0) {互不蕴含};
  \node[font=\footnotesize,black!70,align=center] at (2.9,1.5) {反向失败：\\滑动窗口};
  \node[font=\footnotesize,black!70,align=center] at (2.9,-1.5) {反向失败：\\越来越高的尖峰};
  \node[font=\footnotesize,black!70,align=center] at (8.0,0.75) {反向失败：\\同分布而独立};
\end{tikzpicture}

\emph{箭头都是单向的，每根箭头的反方向都被一个具体反例挡住；左侧两种强收敛之间连箭头都没有，谁也推不出谁。本节余下的篇幅就是把这三个反例造出来。}
\end{center}

几乎必然收敛（almost surely，简写 a.s.）盯的是单条轨迹：\(P(\set{\omega: X_n(\omega)\to X(\omega)})=1\)。固定一个随机种子跑到底，得到的是一列普通实数，它按微积分意义收敛；允许有一个概率为零的例外种子集合，那里可以随便乱来。``几乎''二字不能省——在连续的样本空间上要求每一个 \(\omega\) 都收敛，通常做不到，也没必要。

依概率收敛盯的是每一时刻的失败概率：对每个 \(\varepsilon>0\)，\(P(\abs{X_n-X}>\varepsilon)\to0\)。它从不追问同一个 \(\omega\) 在不同 \(n\) 上的连贯表现，只问``第 \(n\) 步偏离超过 \(\varepsilon\) 的那批样本点占多大分量''。坏样本点可以每一步换一批人。

\(L^p\) 收敛盯的是平均误差：\(\E[\abs{X_n-X}^p]\to0\)。\(p=2\) 就是均方误差趋于零，这是机器学习里最常用的那一种。

依分布收敛只看直方图：在极限分布函数 \(F_X\) 的每个连续点 \(x\) 上有 \(F_{X_n}(x)\to F_X(x)\)。它甚至不要求 \(X_n\) 与 \(X\) 住在同一个概率空间里，两条曲线的形状趋于一致就算数。定义里限定``连续点''是为了排除一个荒唐结论：常数列 \(X_n\equiv1+1/n\) 在跳跃点 \(x=1\) 处的分布函数值恒为 \(0\)，而极限的值是 \(1\)，若不排除跳跃点，连这种最平凡的收敛都会被判为失败。

三根正向箭头都不难。\(L^p\) 推出依概率，只需一行 Markov 不等式：
\[
P(\abs{X_n-X}>\varepsilon)=P(\abs{X_n-X}^p>\varepsilon^p)\le\frac{\E[\abs{X_n-X}^p]}{\varepsilon^p}\to0.
\]
几乎必然推出依概率，是因为整条尾部都待在 \(\varepsilon\) 邻域内这件事，其概率趋于一，那么单看第 \(n\) 项自然也趋于一。依概率推出依分布则需要一段夹逼：把 \(P(X_n\le x)\) 上下用 \(F_X(x\pm\varepsilon)\) 加一个小概率夹住，再让 \(\varepsilon\to0\) 并调用 \(F_X\) 在 \(x\) 处的连续性。这段验算是机械的，跳过它不影响你使用结论。

反方向才是本节真正的内容。

\subsection{滑动窗口：依概率收敛而没有一条路径收敛}

样本空间取 \((0,1]\)，配均匀概率。把区间按组切分：第 \(k\) 组是 \(k\) 个长度为 \(1/k\) 的小区间，从左到右恰好铺满 \((0,1]\)。把这些小区间按组依次排成一列 \(A_1,A_2,A_3,\ldots\)，令 \(X_n=\ind_{A_n}\)。

\begin{center}
\begin{tikzpicture}[x=1cm,y=1cm,font=\footnotesize]
  \fill[blue!30] (0,0) rectangle (8,0.5);
  \draw[black!60] (0,0) rectangle (8,0.5);
  \node[left,black!70] at (-0.25,0.25) {\(k=1\)};
  \fill[blue!30] (0,-1) rectangle (4,-0.5);
  \draw[black!60] (0,-1) rectangle (8,-0.5);
  \draw[black!60] (4,-1) -- (4,-0.5);
  \node[left,black!70] at (-0.25,-0.75) {\(k=2\)};
  \fill[blue!30] (2.667,-2) rectangle (5.333,-1.5);
  \draw[black!60] (0,-2) rectangle (8,-1.5);
  \draw[black!60] (2.667,-2) -- (2.667,-1.5);
  \draw[black!60] (5.333,-2) -- (5.333,-1.5);
  \node[left,black!70] at (-0.25,-1.75) {\(k=3\)};
  \fill[blue!30] (2,-3) rectangle (4,-2.5);
  \draw[black!60] (0,-3) rectangle (8,-2.5);
  \draw[black!60] (2,-3) -- (2,-2.5);
  \draw[black!60] (4,-3) -- (4,-2.5);
  \draw[black!60] (6,-3) -- (6,-2.5);
  \node[left,black!70] at (-0.25,-2.75) {\(k=4\)};
  \draw[dashed,black!75] (3.0,0.85) -- (3.0,-3.3);
  \node[above,black!75] at (3.0,0.85) {固定的 \(\omega\)};
  \node[black!70] at (9.6,0.25) {\(P(A_n)=1\)};
  \node[black!70] at (9.6,-0.75) {\(1/2\)};
  \node[black!70] at (9.6,-1.75) {\(1/3\)};
  \node[black!70] at (9.6,-2.75) {\(1/4\)};
\end{tikzpicture}

\emph{每一组的小区间铺满整个 \((0,1]\)，所以任何固定的 \(\omega\) 在每一组里都恰好被命中一次，涂色格子是它落入的那一个。窗口长度按 \(1/k\) 缩小，坏事件的概率趋于零；但窗口一直在横向滑动，从不放过任何一个 \(\omega\)。}
\end{center}

第 \(n\) 项属于第 \(k(n)\) 组时 \(P(A_n)=1/k(n)\)，而 \(k(n)\) 随 \(n\) 增大到无穷，于是 \(P(\abs{X_n}>1/2)\to0\)，依概率收敛到 \(0\) 成立。

再看路径。图上那条竖虚线是任意一个固定的 \(\omega\)。它在每一组里都被涂色格子命中一次，所以序列 \(X_1(\omega),X_2(\omega),\ldots\) 含有无穷多个 \(1\)，也含有无穷多个 \(0\)，在两个值之间永远弹跳。几乎必然收敛不是``绝大多数路径收敛''，这里是没有一条路径收敛。

这个例子还顺手回答了一个常被混淆的问题：极限号和概率号不能交换。此处 \(\lim_n P(\abs{X_n}>\varepsilon)=0\)，而 \(P(\lim_n X_n=0)=0\)。两种运算顺序给出完全相反的答案，想交换就得另找理由，通常来自单调收敛、控制收敛或一致可积性（见\hyperref[sec:12-Real-Analysis-Support/05-Measure-Integration-and-Conditional-Expectation/06]{``一致可积性阻止期望质量逃逸''}）。

\subsection{高瘦尖峰：概率消失了，期望没有}

同样在 \((0,1]\) 上，令 \(X_n(\omega)=n\) 当 \(0<\omega\le1/n\)，其余取 \(0\)。取值非零的概率是 \(1/n\)，依概率收敛到 \(0\) 毫无问题。可是
\[
\E[X_n]=n\cdot\frac1n=1
\]
对每个 \(n\) 都等于 \(1\)。更狠一点，\(\E[\abs{X_n}^p]=n^{p-1}\)，\(p>1\) 时直接发散。

原因很朴素：概率缩小的速度恰好被幅度增长的速度抵消。这个序列同时还是几乎必然收敛的——对固定的 \(\omega\)，只要 \(n>1/\omega\) 就永远取零。于是它一口气证明了两件事：依概率不蕴含 \(L^p\)，几乎必然也不蕴含 \(L^p\)。图上左侧两个框之间那条``互不蕴含''，一半靠它。另一半靠一列振幅递减而振动不止的函数：均方误差趋于零，可是每个点上都在抖，路径不收敛。

对做机器学习的人，这条裂缝有一个直接后果。``估计量以高概率接近真值''和``估计量的均方误差小''是两笔账。一个罕见但灾难性的大偏差不会显著抬高前者，却能独自撑起后者。如果你的损失是平方型的，报告依概率收敛就不够。

\subsection{依分布收敛最弱：形状一样，位置可以完全无关}

反例简单到近乎无赖。取一列 \(X_n\) 与 \(X\) 同分布，且与 \(X\) 独立。分布函数逐点相同，依分布收敛平凡成立。但 \(X_n-X\) 的分布不随 \(n\) 变窄，除非共同分布本身退化成常数。依分布收敛管形状，不管两个变量是否站在一起。

有一个重要例外：极限是常数 \(c\) 时，依分布收敛与依概率收敛等价。常数的分布函数只在 \(c\) 处跳一次，所以 \(c\pm\varepsilon\) 都是连续点，依分布收敛给出 \(P(X_n\le c-\varepsilon)\to0\) 与 \(P(X_n\le c+\varepsilon)\to1\)，两条合起来就是依概率收敛。这个等价是 Slutsky 定理的地基，也解释了为什么``估计量一致''与``估计量依分布收敛到真值''在实践中被当作同一句话。

还有一条方向相反的桥：若 \(X_n\) 依概率收敛到 \(X\)，则存在子列几乎必然收敛。选 \(n_k\) 使 \(P(\abs{X_{n_k}-X}>2^{-k})<2^{-k}\)，概率和收敛，Borel--Cantelli 引理断言坏事件只发生有限多次。需要把``大概率接近''升级成``路径接近''时，这是标准出口。

\subsection{选哪一种，取决于你要保护什么判断}

分类不是为了考试。四种收敛对应四种可兑现的承诺，用错了名字，定理就保护不到你真正关心的东西。

估计量的一致性用依概率收敛表述。你只需要控制一次实验里判断出错的概率，不需要保证每条轨迹最终安分。极大似然估计、矩估计的一致性定理，结论都停在这一层。

风险的收敛要 \(L^2\)。上面的尖峰反例已经说明，依概率收敛管不住期望，更管不住平方期望。

一条具体轨迹的长期行为要几乎必然收敛。这一条对数据工作者格外要紧：你在一台机器上只跑一条路径，几乎必然收敛承诺的正是``你眼前这一条就是极限''。强大数定律、遍历定理、随机优化与 MCMC 的收敛诊断，都需要这个层级的保证。SGD 的迭代序列是否几乎必然进入驻点邻域，与它是否依概率进入，是两个不同的定理（参见\hyperref[sec:09-Convex-Optimization/08-Nonconvex-and-Stochastic-Optimization/04]{``SGD 在期望中下降但噪声会留下平台''}）。

置信区间与假设检验只需要依分布收敛。这是四种里最弱的一种，也是应用最广的一种——你多半并不需要关于某次运行的承诺，你需要的是抽样分布的形状。

最后补一句容易被忽略的：如果最终目的是让 \(\E[X_n]\to\E[X]\)，上面任何一种收敛单独都不够，还得加一致可积性。这正是尖峰反例踩中的地方。

回到开头那句``收敛了''。下一节讨论样本均值趋于期望，它有两个版本：弱大数定律给的是依概率，强大数定律给的是几乎必然。看完本节这张图你就知道，这两句话的差距在哪里，以及为什么需要多花好几倍的力气去证后一句。
